%% kapa: exam=1 round=23 year=2012 period=1 form=A subject=economics subject_ko=경제원론 seq=28 number=68
%% source: Q-Net 감정평가사 기출문제 (cst00309) artlSeq=5052578 file=raw/1008/5052578/2012년 제23회 감정평가사 제1차시험 기출문제(1교시 A형).pdf
%% conversion: kordoc 4.24.51 -> tools/gen1.py (auto); review=proofread-sonnet; flags=none
\begin{question}{68}
인플레이션 조세(inflation tax)에 관한 설명으로 옳은 것은?
\begin{choices}
  \item 물가가 상승함에 따라 납세자들이 더 높은 세율등급을 적용받아 납부하는 소득세로 정의된다.
  \item 물가가 상승함에 따라 경제주체가 보유하고 있는 통화의 실질가치가 상승할 때 발생한다.
  \item 세율이 인상됨에 따라 인플레이션율이 상승하는 것을 의미한다.
  \item 정부가 정부채권을 시중금융기관으로부터 매입함으로써 발생한 이자율 하락으로 인한 금융자산의 가격하락을 의미한다.
  \item 정부가 통화량을 증가시켜 재정자금을 조달할 때 발생한다.
\end{choices}
\end{question}
\answer{5}
